\par\Needspace{6\baselineskip}\medskip
\begin{center}\hypertarget{native-197-article147}{}\pdfbookmark[2]{147.}{article-147}147.\end{center}
Strictly speaking, here as in the ten methods presented above, we could choose at will what new values of $x$ and $y$ to assume in the second hypothesis, provided they satisfy the general conditions explained above. Yet being able to start at once from more accurate values is clearly a great advantage, so it would be unwise in this method to choose the second values at random. From the very nature of the problem it is easy to see that if the first values of $P,Q$ are affected by small errors, $P',Q'$ will give much more accurate values, provided the heliocentric motion is moderate. We shall therefore always adopt $P',Q'$ as the second values of $P,Q$, or $\log P',\log Q'$ as the second values of $x,y$, if $\log P,\log Q$ were assumed to designate the first.

In this second hypothesis all the preliminary operations given by formulas 1—20 are retained unchanged, and the calculation is repeated in exactly the same way. First determine the angle $\omega$, then $z,r',\dfrac{n'r'}{n}\mathpunct{\raisebox{.45ex}{,}}$ $\dfrac{n'r'}{n''}$
$\zeta,r,\zeta'',r'',f',f,f''$. From the greater or smaller difference between the new values of these quantities and the first ones it is easy to judge whether it is worth recalculating the correction of the times for aberration as well. If it is not, the time intervals and hence the quantities $\theta$ and $\theta''$ remain as before. Finally, from $f,r',r''$; $f'',r,r'$ and the time intervals, obtain $\eta,\eta''$ and hence new values of $P',Q'$. These will usually differ from those given by the first hypothesis much less than those in turn differ from the first values of $P,Q$. The second values of $X,Y$ will thus be much
\SourceLeaf{198}smaller than the first. The second values of $P',Q'$ are adopted as the third values of $P,Q$, and the calculation is repeated with them.
Thus just as the second hypothesis yielded more accurate numbers than the first, the third will in turn yield more accurate ones than the second. The third values of $P',Q'$ could be adopted as the fourth values of $P,Q$, and the calculations repeated until a hypothesis is reached in which $X$ and $Y$ can be treated as vanishing. But whenever the third hypothesis still seems insufficient, it is preferable to derive the values of $P,Q$ to be adopted in the fourth hypothesis from the first three by the method explained in Arts. 120, 121. In this way the approximation is obtained more rapidly, and there will seldom be any need to proceed to a fifth hypothesis.
